\documentclass[11pt,respuestas,a4]{aleph-examen}
% \documentclass[10pt,a5]{aleph-examen}

% -- Paquetes extra
\usepackage{aleph-comandos}
\usepackage{systeme}

% -- Datos
\institucion{Facultad de Ciencias Exactas, Naturales y Ambientales}
\carrera{Catálogo STEM}
\asignatura{Matemática III}
\tema{Examen Escrito 2 -- Criterio 3.2}
\autor{Andrés Merino}
\fecha{Periodo 2026-1}

\logouno[0.16\textwidth]{Logos/logoPUCE_04_ac}
\definecolor{colortext}{HTML}{1957d1}
\definecolor{colordef}{HTML}{1957d1}
\fuente{montserrat}


\begin{document}

\encabezado

\vspace{-5mm}
\section*{Indicaciones}
\begin{itemize}[leftmargin=*]
\item
    En esta actividad se evalúa si el estudiante (Criterio 3.2) emplea la derivada para resolver problemas de optimización.
\item
    Los ejercicios deben ser resueltos a mano y entregarse en hojas de forma física al docente, sin ayuda de resúmenes o asistentes computacionales.
\item
    Se encuentra prohibido el uso de cualquier otra fuente de información durante todo el examen.
\item
    No está permitido tener ningún dispositivo electrónico encendido durante el examen, incluyendo celulares, tablets, computadoras, audífonos, relojes inteligentes, ni similares.
\item
    En caso de considerar que existe un error en la pregunta o que esta se encuentra mal planteada, se debe indicar cuál es el error y justificarlo.
\item
    Todas las soluciones deben estar correctamente redactadas y explicadas.
\end{itemize}

\section*{Ejercicios}

\begin{preguntas}

%%%%%%%%%%%%%%%%%%%%%%%%%%%%%%%%%%%%%%%%%%
\item
Encuentre y clasifique los puntos críticos de la función\puntaje{10}
\[
    f(x,y) = 2xy - 2x^2 - 5y^2 + 4y - 3.
\]

\begin{respuesta}
Iniciemos encontrando los puntos críticos. Para ello, calculemos las derivadas parciales de $f$:
\[
    \frac{\partial f}{\partial x}(x,y) = 2y - 4x,
    \qquad
    \frac{\partial f}{\partial y}(x,y) = 2x - 10y + 4.
\]
Igualando a cero, obtenemos el sistema:
\[
    \left\{
    \begin{aligned}
        2y - 4x &= 0,\\
        2x - 10y + 4 &= 0.
    \end{aligned}
    \right.
\]
De la primera ecuación, $y = 2x$. Sustituyendo en la segunda:
\[
    2x - 10(2x) + 4 = 0
    \;\Longrightarrow\;
    -18x + 4 = 0
    \;\Longrightarrow\;
    x = \frac{2}{9},
    \qquad
    y = \frac{4}{9}.
\]
Con esto, el único punto crítico es $\left(\dfrac{2}{9},\dfrac{4}{9}\right)$.

Apliquemos ahora el criterio de la segunda derivada. Calculemos la matriz Hessiana:
\[
    H_f(x,y)
    =
    \begin{pmatrix}
        \dfrac{\partial^2 f}{\partial x^2} & \dfrac{\partial^2 f}{\partial x\partial y}\\[10pt]
        \dfrac{\partial^2 f}{\partial y\partial x} & \dfrac{\partial^2 f}{\partial y^2}
    \end{pmatrix}
    =
    \begin{pmatrix}
        -4 & 2\\
        2 & -10
    \end{pmatrix}.
\]
La Hessiana es constante; aplicando el criterio de Sylvester:
\[
    \det(H_f) = (-4)(-10) - (2)(2) = 36 > 0,
    \qquad
    \frac{\partial^2 f}{\partial x^2} = -4 < 0,
\]
por lo que la función alcanza un {máximo} local en $\left(\dfrac{2}{9},\dfrac{4}{9}\right)$ y el máximo es
\[
    f\!\left(\frac{2}{9},\frac{4}{9}\right) = -\frac{19}{9} \approx -2{,}11.\qedhere
\]
\end{respuesta}


%%%%%%%%%%%%%%%%%%%%%%%%%%%%%%%%%%%%%%%%%%
\item
Utilice el método de multiplicadores de Lagrange para encontrar los valores extremos de
\[
    f(x,y) = 5x + 2y
\]
sujeta a la restricción $5x^2 + 2y^2 = 14$.\puntaje{15}

\begin{respuesta}
Usemos multiplicadores de Lagrange con $g(x,y) = 5x^2 + 2y^2 - 14$. Planteemos el Lagrangeano:
\[
    L(x,y,\lambda) = 5x + 2y - \lambda(5x^2 + 2y^2 - 14).
\]
Ahora, encontremos los puntos críticos. Para ello, calculemos las derivadas parciales de $L$:
\[
    \frac{\partial L}{\partial x} = 5 - 10\lambda x,
    \qquad
    \frac{\partial L}{\partial y} = 2 - 4\lambda y,
    \qquad
    \frac{\partial L}{\partial \lambda} = -(5x^2 + 2y^2 - 14).
\]
Igualando a cero, obtenemos el sistema:
\[
    \left\{
    \begin{aligned}
        5 - 10\lambda x &= 0,\\
        2 - 4\lambda y &= 0,\\
        5x^2 + 2y^2 - 14 &= 0.
    \end{aligned}
    \right.
\]
De las dos primeras ecuaciones, $x = \dfrac{1}{2\lambda}$ e $y = \dfrac{1}{2\lambda}$, es decir, $x = y$. Sustituyendo en la tercera:
\[
    5x^2 + 2x^2 - 14 = 0
    \;\Longrightarrow\;
    7x^2 = 14
    \;\Longrightarrow\;
    x^2 = 2
    \;\Longrightarrow\;
    x = \pm\sqrt{2}.
\]
Con esto, los puntos críticos son $\bigl(\sqrt{2},\sqrt{2}\bigr)$ con $\lambda = \dfrac{\sqrt{2}}{4}$ y $\bigl(-\sqrt{2},-\sqrt{2}\bigr)$ con $\lambda = -\dfrac{\sqrt{2}}{4}$.

Apliquemos ahora el criterio de la segunda derivada. Calculemos la matriz Hessiana de $L$ con respecto a $(\lambda, x, y)$:
\[
    H_L(\lambda,x,y)
    =
    \begin{pmatrix}
        0 & -10x & -4y\\[4pt]
        -10x & -10\lambda & 0\\[4pt]
        -4y & 0 & -4\lambda
    \end{pmatrix}.
\]
Un cálculo directo muestra que $\det(H_L) = \lambda\bigl(400x^2 + 160y^2\bigr)$. Evaluando en cada punto crítico:
\begin{itemize}
    \item {$\bigl(\sqrt{2},\sqrt{2}\bigr)$ con $\lambda = \dfrac{\sqrt{2}}{4}$:} se tiene $\det(H_L) = 280\sqrt{2} > 0$, por lo que $f$ alcanza un {máximo} y el máximo es
    \[
        f\bigl(\sqrt{2},\sqrt{2}\bigr) = 5\sqrt{2} + 2\sqrt{2} = 7\sqrt{2} \approx 9{,}90.
    \]
    \item {$\bigl(-\sqrt{2},-\sqrt{2}\bigr)$ con $\lambda = -\dfrac{\sqrt{2}}{4}$:} se tiene $\det(H_L) = -280\sqrt{2} < 0$, por lo que $f$ alcanza un {mínimo} y el mínimo es
    \[
        f\bigl(-\sqrt{2},-\sqrt{2}\bigr) = -7\sqrt{2} \approx -9{,}90.\qedhere
    \]
\end{itemize}
\end{respuesta}


%%%%%%%%%%%%%%%%%%%%%%%%%%%%%%%%%%%%%%%%%%
\item
Se va a fabricar una lata cilíndrica de metal (con tapa y fondo) a partir de $54\pi\ \text{cm}^2$ de lámina de lata. ¿Cuál es la mayor capacidad que puede tener la lata? {\footnotesize(\noindent\textit{Pista:} el volumen de un cilindro  es $\pi r^2 h$ y su área total (con tapa y fondo) es $2\pi r^2 + 2\pi r h$)}\puntaje{15}

\begin{respuesta}
Sean $r$ el radio de la base (en cm) y $h$ la altura de la lata (en cm). El volumen es
\[
    V(r,h) = \pi r^2 h,
\]
y, como la lata tiene tapa y fondo, el área total de lámina es
\[
    2\pi r^2 + 2\pi r h = 54\pi.
\]
Así, el problema de optimización es:
\[
    \left\{
    \begin{aligned}
        &\text{max } V(r,h) = \pi r^2 h,\\
        &\text{sujeto a} \\
        &\qquad 2\pi r^2 + 2\pi r h - 54\pi = 0.
    \end{aligned}
    \right.
\]

Usemos multiplicadores de Lagrange. Planteemos el Lagrangeano:
\[
    L(\lambda,r,h) = \pi r^2 h - \lambda\bigl(2\pi r^2 + 2\pi r h - 54\pi\bigr).
\]
Ahora, encontremos los puntos críticos. Para ello, calculemos las derivadas parciales de $L$:
\[
    \frac{\partial L}{\partial r} = 2\pi r h - \lambda(4\pi r + 2\pi h),\quad
    \frac{\partial L}{\partial h} = \pi r^2 - \lambda(2\pi r),\quad
    \frac{\partial L}{\partial \lambda} = -\bigl(2\pi r^2 + 2\pi r h - 54\pi\bigr).
\]
Igualando a cero, obtenemos el sistema:
\[
    \left\{
    \begin{aligned}
        2\pi r h - \lambda(4\pi r + 2\pi h) &= 0,\\
        \pi r^2 - 2\pi\lambda r &= 0,\\
        2\pi r^2 + 2\pi r h - 54\pi &= 0.
    \end{aligned}
    \right.
\]
De la segunda ecuación, $r = 2\lambda$, es decir, $\lambda = \dfrac{r}{2}$. Sustituyendo en la primera:
\[
    2\pi r h - \frac{r}{2}(4\pi r + 2\pi h)  = 0 
    \;\Longrightarrow\;
    2\pi r^2 + \pi r h = 0
    \;\Longrightarrow\;
    \pi r h = 2\pi r^2
    \;\Longrightarrow\;
    h = 2r.
\]
Reemplazando en la tercera:
\[
    2\pi r^2 + 2\pi r(2r) - 54\pi = 0
    \;\Longrightarrow\;
    6\pi r^2 = 54\pi
    \;\Longrightarrow\;
    r^2 = 9
    \;\Longrightarrow\;
    r = 3,\qquad h = 6.
\]
Con esto, el único punto crítico es $(r,h) = (3,6)$ con $\lambda = \dfrac{3}{2}$.

Apliquemos ahora el criterio de la segunda derivada. La matriz Hessiana de $L$ con respecto a $(\lambda,r,h)$, evaluada en el punto crítico, es
\[
    H_L\!\left(\frac{3}{2},3,6\right)
    =
    \begin{pmatrix}
        0 & -24\pi & -6\pi\\[4pt]
        -24\pi & 6\pi & 3\pi\\[4pt]
        -6\pi & 3\pi & 0
    \end{pmatrix},
\]
cuyo determinante es $\det(H_L) = 648\pi^3 > 0$, por lo que $V$ alcanza un {máximo} en $(3,6)$. La mayor capacidad de la lata es
\[
    V(3,6) = \pi (3)^2 (6) = 54\pi \approx 169{,}65\ \text{cm}^3.\qedhere
\]
\end{respuesta}


%%%%%%%%%%%%%%%%%%%%%%%%%%%%%%%%%%%%%%%%%%
\item
Suponga que, en un problema de optimización sin restricciones, una función $f$ tiene el punto crítico $(1,-2,3)$ y que su matriz Hessiana es
\[
    H_f(x,y,z)
    =
    \begin{pmatrix}
        2x & y+1 & z-3\\
        y+1 & -y & x-2\\
        z-3 & x-2 & z-1
    \end{pmatrix}.
\]
Determine la naturaleza del punto crítico:
\begin{enumerate}[leftmargin=*]
    \item usando los valores propios de la Hessiana;\puntaje{6}
    \item usando el criterio de Sylvester.\puntaje{4}
\end{enumerate}

\begin{respuesta}
Primero, evaluemos la matriz Hessiana en el punto crítico $(1,-2,3)$:
\[
    H_f(1,-2,3)
    =
    \begin{pmatrix}
        2(1) & (-2)+1 & 3-3\\
        (-2)+1 & -(-2) & 1-2\\
        3-3 & 1-2 & 3-1
    \end{pmatrix}
    =
    \begin{pmatrix}
        2 & -1 & 0\\
        -1 & 2 & -1\\
        0 & -1 & 2
    \end{pmatrix}.
\]
\begin{enumerate}
    \item \textbf{Valores propios.} Calculemos el polinomio característico $\det(H_f - \lambda I)$:
    \[
        \det
        \begin{pmatrix}
            2-\lambda & -1 & 0\\
            -1 & 2-\lambda & -1\\
            0 & -1 & 2-\lambda
        \end{pmatrix}
        =
        (2-\lambda)\bigl[(2-\lambda)^2 - 1\bigr] - (2-\lambda)
        =
        (2-\lambda)\bigl[(2-\lambda)^2 - 2\bigr].
    \]
    Igualando a cero, se tiene $2-\lambda = 0$ o $(2-\lambda)^2 = 2$, de donde los valores propios son
    \[
        \lambda_1 = 2-\sqrt{2}\approx 0{,}59,
        \qquad
        \lambda_2 = 2,
        \qquad
        \lambda_3 = 2+\sqrt{2}\approx 3{,}41.
    \]
    Como los tres valores propios son positivos, la Hessiana es definida positiva y, por lo tanto, $f$ alcanza un {mínimo} local en $(1,-2,3)$.

    \item \textbf{Criterio de Sylvester.} Calculemos los menores principales líderes:
    \[
        D_1 = 2 > 0,
        \qquad
        D_2 = \det\begin{pmatrix} 2 & -1\\ -1 & 2 \end{pmatrix} = 3 > 0,
        \qquad
        D_3 = \det(H_f) = 4 > 0.
    \]
    Como $D_1, D_2, D_3$ son todos positivos, la Hessiana es definida positiva y, por lo tanto, $f$ alcanza un {mínimo} local en $(1,-2,3)$, en concordancia con el inciso anterior.\qedhere
\end{enumerate}
\end{respuesta}

\end{preguntas}

\end{document}
